% ============================================================
% Question Bank: ch09-03 Experimental Probability
% Topic Title: Experimental Probability
% CCSS: 7.SP.C.7
% Questions: 30 (18 MC + 7 SA + 5 Graphical)
% ============================================================

% --- Multiple Choice Questions (18) ---

\begin{practiceQuestion}{9-3-q01}{mc}
\begin{questionText}
A coin is flipped $200$ times. Heads comes up $112$ times. What is the experimental probability of heads?
\end{questionText}
\choiceA{$0.50$}
\choiceB{$0.56$}
\choiceC{$0.44$}
\choiceD{$0.88$}
\correctAnswer{B}
\explanation{$P = \frac{112}{200} = 0.56$.}
\end{practiceQuestion}

\begin{practiceQuestion}{9-3-q02}{mc}
\begin{questionText}
A die is rolled $60$ times. A $2$ comes up $8$ times. What is the experimental probability of rolling a $2$?
\end{questionText}
\choiceA{$\frac{2}{60}$}
\choiceB{$\frac{1}{6}$}
\choiceC{$\frac{8}{60}$}
\choiceD{$\frac{8}{52}$}
\correctAnswer{C}
\explanation{$P = \frac{\text{times event occurred}}{\text{total trials}} = \frac{8}{60}$.}
\end{practiceQuestion}

\begin{practiceQuestion}{9-3-q03}{mc}
\begin{questionText}
A spinner lands on blue $18$ times out of $90$ spins. What is the experimental probability of landing on blue?
\end{questionText}
\choiceA{$\frac{1}{5}$}
\choiceB{$\frac{1}{4}$}
\choiceC{$\frac{18}{90}$}
\choiceD{Both A and C}
\correctAnswer{D}
\explanation{$P = \frac{18}{90} = \frac{1}{5} = 0.2$. Both A and C represent the same value.}
\end{practiceQuestion}

\begin{practiceQuestion}{9-3-q04}{mc}
\begin{questionText}
Maria flipped a coin $50$ times and got heads $30$ times. The theoretical probability of heads is $0.50$. Which statement is correct?
\end{questionText}
\choiceA{The coin is definitely unfair}
\choiceB{The experimental probability ($0.60$) is higher than the theoretical probability ($0.50$), but this can happen by chance}
\choiceC{The experimental probability always equals the theoretical probability}
\choiceD{Maria made an error because the result should be exactly $25$ heads}
\correctAnswer{B}
\explanation{$P = \frac{30}{50} = 0.60$. Small-sample variation is normal — the experimental probability may differ from the theoretical.}
\end{practiceQuestion}

\begin{practiceQuestion}{9-3-q05}{mc}
\begin{questionText}
A basketball player makes $45$ out of $60$ free throws. What is the experimental probability that she makes her next free throw?
\end{questionText}
\choiceA{$0.45$}
\choiceB{$0.60$}
\choiceC{$0.75$}
\choiceD{$0.25$}
\correctAnswer{C}
\explanation{$P = \frac{45}{60} = 0.75$.}
\end{practiceQuestion}

\begin{practiceQuestion}{9-3-q06}{mc}
\begin{questionText}
A bag contains marbles of unknown colors. After $40$ draws (with replacement), red is drawn $14$ times. What is the experimental probability of drawing red?
\end{questionText}
\choiceA{$\frac{14}{40}$}
\choiceB{$\frac{14}{26}$}
\choiceC{$\frac{26}{40}$}
\choiceD{$\frac{40}{14}$}
\correctAnswer{A}
\explanation{$P = \frac{\text{red draws}}{\text{total draws}} = \frac{14}{40} = 0.35$.}
\end{practiceQuestion}

\begin{practiceQuestion}{9-3-q07}{mc}
\begin{questionText}
As the number of trials in an experiment increases, what happens to the experimental probability?
\end{questionText}
\choiceA{It always equals the theoretical probability}
\choiceB{It gets farther from the theoretical probability}
\choiceC{It tends to get closer to the theoretical probability}
\choiceD{It stays the same regardless of the number of trials}
\correctAnswer{C}
\explanation{The Law of Large Numbers states that as trials increase, experimental probability approaches theoretical probability.}
\end{practiceQuestion}

\begin{practiceQuestion}{9-3-q08}{mc}
\begin{questionText}
A student rolls a die $30$ times and records the results. The number $5$ appears $3$ times. What is the experimental probability of rolling a $5$?
\end{questionText}
\choiceA{$\frac{5}{30}$}
\choiceB{$\frac{3}{30}$}
\choiceC{$\frac{1}{6}$}
\choiceD{$\frac{3}{27}$}
\correctAnswer{B}
\explanation{$P = \frac{3}{30} = \frac{1}{10} = 0.1$.}
\end{practiceQuestion}

\begin{practiceQuestion}{9-3-q09}{mc}
\begin{questionText}
A spinner with $4$ equal sections is spun $80$ times. Section A appears $25$ times. How does the experimental probability compare to the theoretical probability?
\end{questionText}
\choiceA{Experimental ($0.3125$) is higher than theoretical ($0.25$)}
\choiceB{Experimental ($0.25$) equals theoretical ($0.25$)}
\choiceC{Experimental ($0.20$) is lower than theoretical ($0.25$)}
\choiceD{Experimental ($0.3125$) is lower than theoretical ($0.50$)}
\correctAnswer{A}
\explanation{Experimental: $\frac{25}{80} = 0.3125$. Theoretical: $\frac{1}{4} = 0.25$. The experimental is slightly higher.}
\end{practiceQuestion}

\begin{practiceQuestion}{9-3-q10}{mc}
\begin{questionText}
In $100$ spins of a spinner, red appears $22$ times, blue appears $48$ times, and green appears $30$ times. Which color has an experimental probability closest to $\frac{1}{2}$?
\end{questionText}
\choiceA{Red}
\choiceB{Blue}
\choiceC{Green}
\choiceD{Red and Green are equally close}
\correctAnswer{B}
\explanation{Blue: $\frac{48}{100} = 0.48$, which is closest to $0.50$.}
\end{practiceQuestion}

\begin{practiceQuestion}{9-3-q11}{mc}
\begin{questionText}
Tyler draws a card from a standard deck, records the suit, and replaces it. He does this $52$ times. Hearts appear $16$ times. What is the experimental probability of drawing a heart?
\end{questionText}
\choiceA{$\frac{13}{52}$}
\choiceB{$\frac{16}{52}$}
\choiceC{$\frac{16}{36}$}
\choiceD{$\frac{4}{13}$}
\correctAnswer{B}
\explanation{$P = \frac{16}{52} = \frac{4}{13} \approx 0.308$. Note: $\frac{16}{52}$ and $\frac{4}{13}$ are equivalent, but B directly shows the experimental data.}
\end{practiceQuestion}

\begin{practiceQuestion}{9-3-q12}{mc}
\begin{questionText}
A factory tests $500$ lightbulbs and finds $15$ are defective. What is the experimental probability that a lightbulb is defective?
\end{questionText}
\choiceA{$0.015$}
\choiceB{$0.03$}
\choiceC{$0.15$}
\choiceD{$0.30$}
\correctAnswer{B}
\explanation{$P = \frac{15}{500} = 0.03$.}
\end{practiceQuestion}

\begin{practiceQuestion}{9-3-q13}{mc}
\begin{questionText}
A student tosses a thumbtack $100$ times. It lands point-up $62$ times. What is the best estimate for the probability of landing point-up?
\end{questionText}
\choiceA{$0.50$}
\choiceB{$0.38$}
\choiceC{$0.62$}
\choiceD{$1.00$}
\correctAnswer{C}
\explanation{$P = \frac{62}{100} = 0.62$. Experimental data is our best estimate.}
\end{practiceQuestion}

\begin{practiceQuestion}{9-3-q14}{mc}
\begin{questionText}
A die is rolled $120$ times. Each number appears the following number of times: $1$: $22$, $2$: $18$, $3$: $20$, $4$: $21$, $5$: $19$, $6$: $20$. Which number had the highest experimental probability?
\end{questionText}
\choiceA{$1$}
\choiceB{$3$}
\choiceC{$4$}
\choiceD{$6$}
\correctAnswer{A}
\explanation{The number $1$ appeared $22$ times, the most of any number. $P(1) = \frac{22}{120} \approx 0.183$.}
\end{practiceQuestion}

\begin{practiceQuestion}{9-3-q15}{mc}
\begin{questionText}
Aisha flips a coin $10$ times and gets $7$ heads. She says the probability of heads is $0.70$. What should she do to get a more reliable estimate?
\end{questionText}
\choiceA{Accept $0.70$ as the true probability}
\choiceB{Flip the coin many more times}
\choiceC{Use a different coin}
\choiceD{Flip the coin fewer times}
\correctAnswer{B}
\explanation{More trials give a more reliable experimental probability. $10$ trials is a small sample.}
\end{practiceQuestion}

\begin{practiceQuestion}{9-3-q16}{mc}
\begin{questionText}
A weather station recorded rain on $36$ out of $120$ days last spring. Based on this data, what is the experimental probability of rain on any spring day?
\end{questionText}
\choiceA{$0.36$}
\choiceB{$0.30$}
\choiceC{$0.25$}
\choiceD{$0.03$}
\correctAnswer{B}
\explanation{$P = \frac{36}{120} = 0.30$.}
\end{practiceQuestion}

\begin{practiceQuestion}{9-3-q17}{mc}
\begin{questionText}
A spinner has $3$ equal sections. After $90$ spins, Section A appeared $35$ times, Section B appeared $28$ times, and Section C appeared $27$ times. Based on this experiment, is the spinner likely fair?
\end{questionText}
\choiceA{No, because the results are not exactly equal}
\choiceB{Yes, because the results are reasonably close to $30$ each}
\choiceC{No, because Section A appeared too many times}
\choiceD{Yes, because $90$ is divisible by $3$}
\correctAnswer{B}
\explanation{For a fair spinner, each section would appear about $30$ times. The results ($35$, $28$, $27$) are close enough to suggest the spinner is likely fair.}
\end{practiceQuestion}

\begin{practiceQuestion}{9-3-q18}{mc}
\begin{questionText}
A bag has red and blue marbles. Sam draws a marble, records the color, and replaces it $50$ times. Red appears $20$ times. Based on this, what fraction of the marbles in the bag are most likely red?
\end{questionText}
\choiceA{$\frac{1}{2}$}
\choiceB{$\frac{2}{5}$}
\choiceC{$\frac{3}{5}$}
\choiceD{$\frac{1}{5}$}
\correctAnswer{B}
\explanation{$P(\text{red}) = \frac{20}{50} = \frac{2}{5}$. This suggests about $\frac{2}{5}$ of the marbles are red.}
\end{practiceQuestion}

% --- Short Answer Questions (7) ---

\begin{practiceQuestion}{9-3-q19}{sa}
\begin{questionText}
A coin is flipped $250$ times. Tails appears $130$ times. What is the experimental probability of tails? Write your answer as a decimal.
\end{questionText}
\correctAnswer{$0.52$}
\explanation{$P = \frac{130}{250} = 0.52$.}
\end{practiceQuestion}

\begin{practiceQuestion}{9-3-q20}{sa}
\begin{questionText}
A die is rolled $180$ times. The number $4$ appears $25$ times. What is the experimental probability of rolling a $4$? Write your answer as a fraction in simplest form.
\end{questionText}
\correctAnswer{$\frac{5}{36}$}
\explanation{$P = \frac{25}{180} = \frac{5}{36}$.}
\end{practiceQuestion}

\begin{practiceQuestion}{9-3-q21}{sa}
\begin{questionText}
A basketball player attempted $80$ free throws and made $56$. What is the experimental probability that she makes her next free throw? Write your answer as a decimal.
\end{questionText}
\correctAnswer{$0.7$}
\explanation{$P = \frac{56}{80} = 0.7$.}
\end{practiceQuestion}

\begin{practiceQuestion}{9-3-q22}{sa}
\begin{questionText}
In an experiment, a spinner is spun $150$ times. Red appears $45$ times. What is the experimental probability of NOT landing on red? Write your answer as a decimal.
\end{questionText}
\correctAnswer{$0.7$}
\explanation{$P(\text{red}) = \frac{45}{150} = 0.3$. $P(\text{not red}) = 1 - 0.3 = 0.7$.}
\end{practiceQuestion}

\begin{practiceQuestion}{9-3-q23}{sa}
\begin{questionText}
A quality inspector checks $400$ items and finds $12$ defective ones. Based on this, predict how many defective items there would be in a batch of $2{,}000$.
\end{questionText}
\correctAnswer{$60$}
\explanation{$P(\text{defective}) = \frac{12}{400} = 0.03$. In $2{,}000$ items: $0.03 \times 2{,}000 = 60$.}
\end{practiceQuestion}

\begin{practiceQuestion}{9-3-q24}{sa}
\begin{questionText}
A student draws a card from a deck, records whether it is red or black, and replaces it. In $80$ draws, $44$ cards are red. What is the experimental probability of drawing a red card? Write your answer as a decimal.
\end{questionText}
\correctAnswer{$0.55$}
\explanation{$P = \frac{44}{80} = 0.55$.}
\end{practiceQuestion}

\begin{practiceQuestion}{9-3-q25}{sa}
\begin{questionText}
A spinner is spun $40$ times. Green appears $10$ times, blue appears $18$ times, and yellow appears $12$ times. What is the experimental probability of landing on blue or yellow? Write your answer as a decimal.
\end{questionText}
\correctAnswer{$0.75$}
\explanation{Blue $+$ Yellow $= 18 + 12 = 30$. $P = \frac{30}{40} = 0.75$.}
\end{practiceQuestion}

% --- Graphical Questions (3 GMC + 2 GSA) ---

\begin{practiceQuestion}{9-3-q26}{gmc}
\begin{questionText}
The bar graph shows the results of rolling a die $60$ times. Based on the graph, what is the experimental probability of rolling a $3$?

\begin{center}
\begin{tikzpicture}[scale=0.55]
  \draw[thick,->] (0,0) -- (0,14) node[above,font=\small] {Frequency};
  \draw[thick,->] (0,0) -- (8,0);
  \foreach \y in {2,4,...,12} {
    \draw[gray!30] (0,\y) -- (7.5,\y);
    \node[left,font=\small] at (0,\y) {\y};
  }
  \fill[funBlue!70] (0.3,0) rectangle (1,8);
  \fill[funBlue!70] (1.3,0) rectangle (2,12);
  \fill[funBlue!70] (2.3,0) rectangle (3,6);
  \fill[funBlue!70] (3.3,0) rectangle (4,10);
  \fill[funBlue!70] (4.3,0) rectangle (5,14);
  \fill[funBlue!70] (5.3,0) rectangle (6,10);
  \node[below,font=\small] at (0.65,0) {1};
  \node[below,font=\small] at (1.65,0) {2};
  \node[below,font=\small] at (2.65,0) {3};
  \node[below,font=\small] at (3.65,0) {4};
  \node[below,font=\small] at (4.65,0) {5};
  \node[below,font=\small] at (5.65,0) {6};
\end{tikzpicture}
\end{center}
\end{questionText}
\choiceA{$\frac{6}{60}$}
\choiceB{$\frac{10}{60}$}
\choiceC{$\frac{12}{60}$}
\choiceD{$\frac{14}{60}$}
\correctAnswer{A}
\explanation{The bar for $3$ reaches $6$. $P = \frac{6}{60} = \frac{1}{10}$.}
\end{practiceQuestion}

\begin{practiceQuestion}{9-3-q27}{gmc}
\begin{questionText}
The frequency table shows the results of spinning a spinner $100$ times. Which color has an experimental probability closest to $\frac{1}{4}$?

\begin{center}
\begin{tabular}{|c|c|}
\hline
\textbf{Color} & \textbf{Frequency} \\
\hline
Red & $30$ \\
\hline
Blue & $24$ \\
\hline
Green & $18$ \\
\hline
Yellow & $28$ \\
\hline
\end{tabular}
\end{center}
\end{questionText}
\choiceA{Red}
\choiceB{Blue}
\choiceC{Green}
\choiceD{Yellow}
\correctAnswer{B}
\explanation{$\frac{1}{4} = 0.25$. Blue: $\frac{24}{100} = 0.24$, which is closest to $0.25$.}
\end{practiceQuestion}

\begin{practiceQuestion}{9-3-q28}{gmc}
\begin{questionText}
The line graph shows the experimental probability of heads after different numbers of coin flips. What happens to the experimental probability as the number of flips increases?

\begin{center}
\begin{tikzpicture}[scale=0.7]
  \draw[thick,->] (0,0) -- (7,0) node[right,font=\small] {Flips};
  \draw[thick,->] (0,0) -- (0,5) node[above,font=\small] {$P$(Heads)};
  \node[left,font=\small] at (0,2.5) {$0.50$};
  \draw[dashed,gray] (0,2.5) -- (6.5,2.5);
  \node[left,font=\small] at (0,0) {$0$};
  \node[left,font=\small] at (0,5) {$1.0$};
  \node[below,font=\small] at (1,0) {10};
  \node[below,font=\small] at (2,0) {20};
  \node[below,font=\small] at (3,0) {50};
  \node[below,font=\small] at (4,0) {100};
  \node[below,font=\small] at (5,0) {200};
  \node[below,font=\small] at (6,0) {500};
  \fill[funRed] (1,3.5) circle (0.08);
  \fill[funRed] (2,3.0) circle (0.08);
  \fill[funRed] (3,2.8) circle (0.08);
  \fill[funRed] (4,2.6) circle (0.08);
  \fill[funRed] (5,2.55) circle (0.08);
  \fill[funRed] (6,2.52) circle (0.08);
  \draw[thick,funRed] (1,3.5) -- (2,3.0) -- (3,2.8) -- (4,2.6) -- (5,2.55) -- (6,2.52);
\end{tikzpicture}
\end{center}
\end{questionText}
\choiceA{It moves away from $0.50$}
\choiceB{It stays exactly at $0.50$}
\choiceC{It gets closer to $0.50$}
\choiceD{It becomes $0$}
\correctAnswer{C}
\explanation{The graph shows that as the number of flips increases, the experimental probability gets closer and closer to the theoretical probability of $0.50$.}
\end{practiceQuestion}

\begin{practiceQuestion}{9-3-q29}{gsa}
\begin{questionText}
The table shows the results of drawing marbles from a bag $50$ times (with replacement). Based on the data, what is the experimental probability of drawing a green marble? Write your answer as a decimal.

\begin{center}
\begin{tabular}{|c|c|}
\hline
\textbf{Color} & \textbf{Times Drawn} \\
\hline
Red & $18$ \\
\hline
Green & $12$ \\
\hline
Blue & $20$ \\
\hline
\end{tabular}
\end{center}
\end{questionText}
\correctAnswer{$0.24$}
\explanation{$P = \frac{12}{50} = 0.24$.}
\end{practiceQuestion}

\begin{practiceQuestion}{9-3-q30}{gsa}
\begin{questionText}
The dot plot shows the outcomes of spinning a spinner $40$ times. Each dot represents one spin. How many more times would you need to spin ``B'' to make its experimental probability equal to $\frac{1}{2}$? Assume the total number of spins stays at $40$.

\begin{center}
\begin{tikzpicture}[scale=0.8]
  \draw[thick] (0,0) -- (5,0);
  \node[below,font=\bfseries] at (1,0) {A};
  \node[below,font=\bfseries] at (2.5,0) {B};
  \node[below,font=\bfseries] at (4,0) {C};
  \draw[thick] (1,-0.1) -- (1,0.1);
  \draw[thick] (2.5,-0.1) -- (2.5,0.1);
  \draw[thick] (4,-0.1) -- (4,0.1);
  % A: 15 dots (3 columns of 5)
  \foreach \r in {1,...,5} { \fill[funRed] (0.85,0.2+\r*0.3) circle (0.1); }
  \foreach \r in {1,...,5} { \fill[funRed] (1.15,0.2+\r*0.3) circle (0.1); }
  \foreach \r in {1,...,5} { \fill[funRed] (1.0,0.2+\r*0.3+1.5) circle (0.1); }
  % B: 12 dots (3 columns of 4)
  \foreach \r in {1,...,4} { \fill[funBlue] (2.35,0.2+\r*0.3) circle (0.1); }
  \foreach \r in {1,...,4} { \fill[funBlue] (2.65,0.2+\r*0.3) circle (0.1); }
  \foreach \r in {1,...,4} { \fill[funBlue] (2.5,0.2+\r*0.3+1.2) circle (0.1); }
  % C: 13 dots (3 cols)
  \foreach \r in {1,...,5} { \fill[funGreen] (3.85,0.2+\r*0.3) circle (0.1); }
  \foreach \r in {1,...,5} { \fill[funGreen] (4.15,0.2+\r*0.3) circle (0.1); }
  \foreach \r in {1,...,3} { \fill[funGreen] (4.0,0.2+\r*0.3+1.5) circle (0.1); }
  \node[above,font=\small] at (1,3.5) {$15$};
  \node[above,font=\small] at (2.5,2.8) {$12$};
  \node[above,font=\small] at (4,3.2) {$13$};
\end{tikzpicture}
\end{center}
\end{questionText}
\correctAnswer{$8$}
\explanation{For $P(B) = \frac{1}{2}$, we need $B = \frac{40}{2} = 20$. Currently $B = 12$, so $20 - 12 = 8$ more spins of B are needed (with $8$ fewer from A and C combined).}
\end{practiceQuestion}
